\documentclass[11pt,a4paper]{article}

\usepackage[utf8]{inputenc}
\usepackage[T1]{fontenc}
\usepackage[french]{babel}
\usepackage{lmodern}
\usepackage[margin=2cm]{geometry}
\usepackage{xcolor}
\usepackage{tikz}
\usepackage{booktabs}
\usepackage{tabularx}
\usepackage{longtable}
\usepackage{array}
\usepackage{siunitx}
\usepackage{amsmath}
\usepackage{amssymb}
\usepackage{pifont}
\usepackage{enumitem}
\usepackage{listings}
\usepackage[colorlinks=true, linkcolor=blue!50!black, urlcolor=blue!60!black]{hyperref}

\usetikzlibrary{arrows.meta, positioning, shapes.geometric, calc, fit, backgrounds}

% Couleurs
\definecolor{chw}{HTML}{B71C1C}
\definecolor{cdsp}{HTML}{1565C0}
\definecolor{ca53}{HTML}{2E7D32}
\definecolor{cusb}{HTML}{6A1B9A}
\definecolor{cnpu}{HTML}{E65100}
\definecolor{cfx}{HTML}{00838F}
\definecolor{cbuf}{HTML}{FFB300}
\definecolor{ccomp}{HTML}{455A64}
\definecolor{ctest}{HTML}{00695C}
\definecolor{codebg}{HTML}{F5F5F5}

% Styles TikZ
\tikzset{
  hw/.style={draw, fill=chw!15, rounded corners=2pt, font=\sffamily\footnotesize, align=center, minimum height=7mm, inner sep=3pt},
  dsp/.style={draw, fill=cdsp!15, rounded corners=2pt, font=\sffamily\footnotesize, align=center, minimum height=7mm, inner sep=3pt},
  a53/.style={draw, fill=ca53!15, rounded corners=2pt, font=\sffamily\footnotesize, align=center, minimum height=7mm, inner sep=3pt},
  comp/.style={draw, fill=ccomp!15, rounded corners=2pt, font=\sffamily\footnotesize, align=center, minimum height=6mm, minimum width=12mm, inner sep=2pt},
  buf/.style={draw, fill=cbuf!25, font=\sffamily\scriptsize, minimum height=5mm, minimum width=6mm, inner sep=1pt, align=center},
  pcm/.style={draw, fill=yellow!30, font=\sffamily\scriptsize, minimum height=5mm, rounded corners=1pt, align=center},
  dai/.style={draw, fill=chw!20, font=\sffamily\footnotesize, align=center, minimum height=8mm, minimum width=14mm, rounded corners=2pt},
  arr/.style={-{Stealth[length=2mm]}, thick},
  arrmulti/.style={-{Stealth[length=2mm]}, line width=1.2pt},
  zone/.style={draw, dashed, rounded corners=6pt, inner sep=8pt},
  pipe/.style={draw, dotted, rounded corners=4pt, inner sep=6pt, fill=cdsp!5}
}

% Code listings
\lstset{
  basicstyle=\ttfamily\footnotesize,
  backgroundcolor=\color{codebg},
  frame=single,
  rulecolor=\color{black!30},
  xleftmargin=0.5em,
  xrightmargin=0.5em,
  breaklines=true,
  columns=flexible,
  keepspaces=true,
  tabsize=2,
  inputencoding=utf8,
  extendedchars=true,
  literate={é}{{\'e}}1 {è}{{\`e}}1 {à}{{\`a}}1 {ê}{{\^e}}1 {â}{{\^a}}1 {ô}{{\^o}}1 {ù}{{\`u}}1 {û}{{\^u}}1 {î}{{\^\i}}1 {ï}{{\"\i}}1 {ç}{{\c c}}1 {É}{{\'E}}1 {È}{{\`E}}1 {À}{{\`A}}1
}

\title{\textbf{Plateforme Mastering Live Debix}\\[6pt]
  \Large Architecture C --- Spécification complète Phase 1a \textbf{V2}\\[4pt]
  \large Structure DSP, tests, livrables et procédure d'exécution}
\author{Document de référence technique}
\date{2026-04-21 --- rev V2}

\begin{document}
\maketitle
\tableofcontents
\clearpage

% =======================================================================
\section*{Changelog V1 $\rightarrow$ V2}
\addcontentsline{toc}{section}{Changelog V1 $\rightarrow$ V2}

Cette V2 intègre les findings d'une investigation critic (6 workers IA, 18 fichiers analysés). Les corrections majeures :

\begin{center}
\begin{tabularx}{\textwidth}{llX}
\toprule
\textbf{Ref} & \textbf{Sévérité} & \textbf{Correction V2} \\
\midrule
B2 & Critique & \texttt{PCM\_DUPLEX\_ADD(TAC5212\_SAI, ...)} $\rightarrow$ remplacé par \texttt{PCM\_CAPTURE\_ADD(SAI\_Capture,...)} + \texttt{PCM\_PLAYBACK\_ADD(SAI\_Playback,...)} séparés. Un device duplex ne peut pas exposer 2 noms distincts. \\
B3 & Critique & Test T4 null-test : remplacé la comparaison \texttt{sox -m} par un script Python avec cross-correlation pour aligner source et capture au sample près (latence pipeline). \\
B4 & Critique & Test T8 : typo \texttt{plychw} $\rightarrow$ \texttt{plughw}. \\
B7 & Critique & PIPE 5 (Master\_Tap) : ajout \texttt{PIPELINE\_SCHED\_COMP\_3} pour que PIPE 3 et PIPE 5 partagent le même scheduler. Empêche le drift / race entre 2 timers indépendants. \\
B11 & Critique & Squelette m4 : placeholders \texttt{...} remplacés par valeurs concrètes (buffer sizes, COMP\_SAMPLE\_SIZE, PIPELINE\_CHANNELS). Fichier désormais compilable. \\
HIGH-2 & Haute & Volume soft-ramp : procédure T4 attend 100 ms après chaque \texttt{amixer cset} pour laisser la rampe converger avant null-test. \\
B8 & Haute & Budget latence mic$\rightarrow$HP : documentation du jitter timer SOF (\textpm 500 µs) + group delay SigmaDelta TAC5212 (1 ms ADC + 0.5 ms DAC). Cible réelle 5-6 ms, pas 4 ms. Critère T8 relâché à $\leq$ 6 ms. \\
B5 / T9 & Moyenne & Seuil CPU T9 : 30\,\% $\rightarrow$ \textbf{35\,\%} pour compter l'overhead des memcpy host bridges. \\
Tooling & Moyenne & Ajout \texttt{sof-tools} dans \texttt{IMAGE\_INSTALL} + déploiement du \texttt{.ldc} à côté du \texttt{.ri} (requis pour T9 \texttt{sof-logger}). \\
Nom carte & Moyenne & Ajout T0 pré-test : \texttt{cat /proc/asound/cards} pour confirmer le nom de la carte ALSA avant d'utiliser \texttt{plughw:softac5212tdm}. \\
Gate V2 & Moyenne & Ajout étape obligatoire « compile-test m4 » avant tout deploy sur la board. \\
\bottomrule
\end{tabularx}
\end{center}

\noindent\textit{Findings considérés comme faux positifs (dus aux submodules \texttt{sof/} non clonés par l'investigation) : la "baseline fantôme" drc2ms.m4 existe bien dans le submodule sof/ branche feature ; le "dummy\_dma CPU memcpy" est résolu par le fork local SOF avec memcpy\_dma + AP2AP fonctionnel.}

\clearpage

% =======================================================================
\section{Rappel : Architecture C}

\subsection{Décision validée}
Architecture retenue par l'utilisateur (voir \texttt{DSP\_ARCHITECTURES.pdf} --- comparatif 3 archis). Architecture C dite "équilibrée" car elle répartit la charge entre les 3 étages (codec TAC5212, DSP HiFi4, Cortex-A53 userspace) selon leurs forces.

\subsection{Principe directeur}
\begin{itemize}[leftmargin=*, itemsep=2pt]
  \item \textbf{Codec TAC5212} (I\textsuperscript{2}C) : pré-conditionnement per-channel --- AGC, HPF, biquads, gain/phase calibration, DVC. \textit{Décharge le DSP du traitement par canal.}
  \item \textbf{DSP HiFi4 (SOF)} : ce qui DOIT être en basse latence --- routage matrice, chaîne master stéréo, tap NPU. \textit{Latence déterministe.}
  \item \textbf{A53 userspace} : USB gadget bridges, JACK+LV2 pour les effets send (reverb, chorus, flanger, delay), service NPU, UI Flutter. \textit{Latence non critique.}
\end{itemize}

\subsection{Vue macro système}

\begin{center}
\begin{tikzpicture}[node distance=5mm and 12mm]

\node[hw]                   (mics)  {8 mics / lines};
\node[hw, below=2mm of mics](tac)   {4× TAC5212\\AGC + 3 biquads/ch\\+ HPF + phase cal\\+ DVC};

\node[dsp, right=18mm of tac, minimum width=60mm, minimum height=55mm] (dsp) {};
\node[anchor=north] at (dsp.north) {\textbf{\textcolor{cdsp}{DSP HiFi4 (SOF, Phase 1a)}}};
\node[dsp, fill=cdsp!25] (eq) at ($(dsp.west)+(1.2,1.8)$) {8× eq\_iir\\(fine EQ)};
\node[dsp, fill=cdsp!25, right=5mm of eq] (matrix) {Routing\\fixe Phase 1a};
\node[dsp, fill=cdsp!25, below=3mm of matrix] (master) {Master chain\\(mbdrc + vol\\+ drc lim)};
\node[dsp, fill=cdsp!25, left=5mm of master] (tap) {NPU tap};
\node[dsp, fill=cdsp!25, below=3mm of master] (mux) {mux 8ch\\(merge master\\+ monitor)};

\node[a53, right=22mm of dsp, minimum width=55mm, minimum height=55mm] (a53) {};
\node[anchor=north] at (a53.north) {\textbf{\textcolor{ca53}{A53 userspace}}};
\node[a53, fill=cfx!25] (fx) at ($(a53.west)+(1.1,1.6)$) {JACK + 4 LV2\\(Phase 3)\\Reverb Chorus\\Flanger Delay};
\node[a53, fill=cusb!25, below=3mm of fx] (usb) {UAC2 gadgets\\(Phase 2)\\8+8 + 2+2};
\node[a53, fill=cnpu!25, below=3mm of usb] (npu) {NPU service\\(Phase 4)};

\node[hw, left=14mm of mux] (hp) {HP / monitors\\(4× TAC5212 DAC)};

\draw[arrmulti] (mics) -- (tac);
\draw[arrmulti] (tac.east) -- node[above, font=\scriptsize, align=center]{SAI7 TDM\\8ch 32b 48k} (eq.west);
\draw[arr] (eq) -- (matrix);
\draw[arr] (matrix.south) -- (master.north);
\draw[arr] (master) -- (mux);
\draw[arr] (master.west) -- (tap.east);
\draw[arrmulti] (mux.west) -- node[above, font=\scriptsize]{SAI7 TX} (hp);
\draw[arr, dashed] (matrix.west) to[bend right=20] (mux);

\draw[arrmulti, <->, dashed] (matrix.east) -- node[above, font=\scriptsize]{(Phase 2--3)} (usb.west);
\draw[arr, dashed, red!70!black] (tap.south) to[bend left=10] (npu.west);
\draw[arr, dashed, red!70!black] (npu.north west) to[bend right=20] node[right, font=\scriptsize, text=red!70!black]{ALSA ctrl} (master.south east);

\end{tikzpicture}
\end{center}
\textit{\small En Phase 1a, seule la partie DSP (bleu) + TAC5212 (rouge) est implémentée. La partie A53 (vert) arrive en Phase 2+.}

\subsection{Contraintes héritées}
\begin{center}
\begin{tabular}{ll}
\toprule
Sample rate & 48 kHz \\
Format & s32le \\
Nb canaux TDM & 8 (TAC5212 $\times$4 daisy-chain) \\
DSP period SOF & 2 ms \\
Latence mic\,$\rightarrow$\,HP direct & $\leq$ 4 ms \\
Latence round-trip USB (plus tard) & $\leq$ 10 ms \\
IPC SOF & IPC3 (hérité platform i.MX8MP) \\
Null-test bypass & bit-perfect requis pour T4 \\
\bottomrule
\end{tabular}
\end{center}

\clearpage
% =======================================================================
\section{Phase 1a --- Objectifs et scope}

\subsection{Objectifs}
\begin{enumerate}[leftmargin=*, itemsep=2pt]
  \item Exposer \textbf{4 PCMs ALSA fonctionnels} sur la carte SOF (\texttt{softac5212tdm}) :
    \begin{center}
    \begin{tabularx}{0.95\textwidth}{llll}
      \toprule
      \textbf{Nom PCM} & \textbf{Type} & \textbf{Ch} & \textbf{Accès userspace typique} \\
      \midrule
      \texttt{SAI\_Capture} & capture & 8 & \texttt{arecord -D plughw:softac5212tdm,SAI\_Capture} \\
      \texttt{SAI\_Playback} & playback & 8 & \texttt{aplay -D plughw:softac5212tdm,SAI\_Playback} \\
      \texttt{Source\_Play} & playback & 2 & \texttt{aplay -D plughw:softac5212tdm,Source\_Play} \\
      \texttt{Master\_Tap} & capture & 2 & \texttt{arecord -D plughw:softac5212tdm,Master\_Tap} \\
      \bottomrule
    \end{tabularx}
    \end{center}
    Ces \textbf{noms} sont explicitement déclarés dans la topology via \texttt{PCM\_DUPLEX\_ADD} / \texttt{PCM\_CAPTURE\_ADD} / \texttt{PCM\_PLAYBACK\_ADD}. On retrouve ces identifiants côté board dans \texttt{arecord -l}, \texttt{/proc/asound/softac5212tdm/}, et dans les scripts de test.
  \item Implanter la \textbf{chaîne master stéréo} complète (multiband\_drc + volume + drc limiter) avec modes bypass.
  \item Garantir \textbf{latence mic\,$\rightarrow$\,HP direct $\leq$ 4 ms}.
  \item Fournir un \textbf{preset bypass bit-perfect} pour le null-test.
  \item Exposer tous les paramètres via \textbf{ALSA controls} pour pilotage live (UI + NPU).
  \item Configurer les \textbf{4 codecs TAC5212} au boot (biquads flat, AGC off, HPF 12 Hz, DVC 0 dB).
  \item Fournir un \textbf{script de test automatisé} couvrant T1\,$\rightarrow$\,T9.
\end{enumerate}

\subsection{Baseline existante --- point de départ}
\textbf{La topology actuellement déployée sur la board fonctionne en duplex}. Rappel :
\begin{itemize}[leftmargin=*, itemsep=1pt]
  \item Fichier actif : \texttt{sof-imx8mp-tac5212-drc2ms.m4} (variante 2 ms)
  \item Déployé en : \texttt{/lib/firmware/imx/sof-tplg/sof-imx8mp-tac5212.tplg}
  \item Validé utilisateur : \texttt{sox -c 8 -t alsa plughw:2,0 -c 8 -t alsa plughw:2,0 remix 8 7 6 5 4 3 2 1} tourne en loopback 8 canaux sans clic ni xrun
  \item Structure : 1 pipeline capture avec DRC + 1 pipeline playback passthrough, exposé en 1 PCM duplex (ID 0)
\end{itemize}
La topology Phase 1a (\texttt{sof-imx8mp-tac5212-masterlite.m4}) \textbf{étend} cette baseline --- elle garde le couple capture/playback SAI (qui marche) et \textbf{ajoute} les pipelines master chain + master tap.

\subsection{Hors scope Phase 1a}
\begin{itemize}[leftmargin=*, itemsep=2pt]
  \item \textcolor{red!70!black}{Matrice N\,$\times$\,M avec gain par cellule} $\rightarrow$ reporté Phase 1b.
  \item USB gadgets UAC2 $\rightarrow$ Phase 2.
  \item Effets send (reverb\ldots) $\rightarrow$ Phase 3.
  \item NPU inférence + training $\rightarrow$ Phase 4.
  \item Harmonizer, Flutter UI avancée $\rightarrow$ Phase 5.
\end{itemize}

\subsection{Stratégie incrémentale d'exécution}
Phase 1a découpée en 3 étapes, \textbf{chacune validée par l'utilisateur avant la suivante} :
\begin{center}
\begin{tabularx}{\textwidth}{lX}
\toprule
\textbf{Étape} & \textbf{Livrables} \\
\midrule
\textbf{1a.1 MVP} & PIPE\_1 (capture dry mics), PIPE\_3+5 (Source\_Play $\rightarrow$ master chain $\rightarrow$ Master\_Tap), PIPE\_6 (SAI\_Playback 8ch direct). Permet tests T1, T2, T3, T4, T6, T7. \\
\textbf{1a.2 Monitor} & Ajout PIPE\_2 (mic\,$\rightarrow$\,HP direct 4 ms). Permet test T8 (latence). \\
\textbf{1a.3 Merge} & Ajout PIPE\_4 via \texttt{mux} pour merge master (2ch, slots TDM 1-2) + monitor (6ch, slots 3-8). Phase 1a complète. Permet tests T5, T9. \\
\bottomrule
\end{tabularx}
\end{center}

\clearpage
% =======================================================================
\section{Structure DSP détaillée}

\subsection{PCMs ALSA exposés}
Côté host A53 (kernel ALSA), la carte \texttt{softac5212tdm} expose 4 devices PCM :

\begin{center}
\begin{tabularx}{\textwidth}{lllX}
\toprule
\textbf{Nom PCM} & \textbf{Type} & \textbf{Canaux} & \textbf{Rôle} \\
\midrule
\texttt{SAI\_Capture} & capture & 8 & Mics bruts post-TAC5212 ADC (avec eq\_iir flat par défaut) \\
\texttt{SAI\_Playback} & playback & 8 & Écriture directe vers TAC5212 DAC (monitoring depuis A53) \\
\texttt{Source\_Play} & playback & 2 & Entrée stéréo master chain (source pour test \texttt{aplay}) \\
\texttt{Master\_Tap} & capture & 2 & Sortie post master chain (NPU + null-test) \\
\bottomrule
\end{tabularx}
\end{center}

Un 5\textsuperscript{e} PCM \texttt{Source\_Play} est ajouté en Phase 1a (vs. le plan initial 3 PCMs) pour permettre le test \texttt{aplay fichier.wav} sans signal mic live.

\subsection{Inventaire des 6 pipelines SOF}

\begin{center}
\begin{tabularx}{\textwidth}{llX}
\toprule
\textbf{Pipeline} & \textbf{Période} & \textbf{Description} \\
\midrule
PIPE\_1 & 2 ms & Capture dry --- SAI7 RX $\rightarrow$ host \texttt{SAI\_Capture} (8ch) \\
PIPE\_2 & 2 ms & Monitor direct mic\,$\rightarrow$\,HP low-latency 4 ms --- SAI7 RX $\rightarrow$ eq\_iir\,$\times$\,8 $\rightarrow$ PIPE\_4 (via buffer) \\
PIPE\_3 & 2 ms & Master chain stéréo --- host \texttt{Source\_Play} $\rightarrow$ multiband\_drc $\rightarrow$ volume $\rightarrow$ drc $\rightarrow$ PIPE\_4 et PIPE\_5 \\
PIPE\_4 & 2 ms & Merge master + monitor --- \texttt{mux} (2ch master sur slots 1-2 + 6ch monitor sur slots 3-8) $\rightarrow$ SAI7 TX (8ch) \\
PIPE\_5 & 2 ms & Master tap host --- PIPE\_3 output (2ch) $\rightarrow$ host \texttt{Master\_Tap} \\
PIPE\_6 & 2 ms & Playback host direct --- host \texttt{SAI\_Playback} (8ch) $\rightarrow$ SAI7 TX (via \texttt{mux} slot override ou pipeline alternatif) \\
\bottomrule
\end{tabularx}
\end{center}

\textit{\small PIPE\_6 est utilisé en 1a.1 (MVP) comme alternative simple à PIPE\_4, puis remplacé/coexistant en 1a.3.}

\subsection{Schéma détaillé des pipelines (Phase 1a complète)}

\begin{center}
\begin{tikzpicture}[node distance=3mm and 7mm]

% PIPE 1 - dry capture
\node[dai] (saiRx) {SAI7 RX\\8ch TDM};
\node[buf, right=6mm of saiRx] (b1) {buf};
\node[pcm, right=6mm of b1] (pcmCap) {SAI\_Capture\\cap 8ch};

\node[pipe, fit=(saiRx) (b1) (pcmCap), label={[anchor=north west, font=\sffamily\bfseries\small, inner sep=2pt]north west:PIPE\_1}] {};

% PIPE 2 - low-latency monitor
\node[buf, below=18mm of saiRx] (b2a) {buf};
\node[comp, right=4mm of b2a] (eqm) {eq\_iir×8};
\node[buf, right=4mm of eqm] (b2b) {buf};

% PIPE 3 - master chain
\node[pcm, below=20mm of b2a] (pcmSrc) {Source\_Play\\play 2ch};
\node[buf, right=4mm of pcmSrc] (b3a) {buf};
\node[comp, right=4mm of b3a, fill=cdsp!30] (mbdrc) {multiband\_drc};
\node[buf, right=4mm of mbdrc] (b3b) {buf};
\node[comp, right=4mm of b3b, fill=cdsp!30] (vol) {volume};
\node[buf, right=4mm of vol] (b3c) {buf};
\node[comp, right=4mm of b3c, fill=cdsp!30] (drclim) {drc (lim)};
\node[buf, right=4mm of drclim] (b3d) {buf};

% PIPE 5 - master tap
\node[pcm, below=10mm of b3d] (pcmTap) {Master\_Tap\\cap 2ch};

% PIPE 4 - mux merge
\node[comp, right=6mm of b3d, fill=cdsp!30, minimum width=18mm, minimum height=20mm] (muxc) {mux\\2ch master\\$\oplus$\\6ch monitor};
\node[buf, right=4mm of muxc] (b4) {buf};
\node[dai, right=4mm of b4] (saiTx) {SAI7 TX\\8ch TDM};

% PIPE 6 - direct playback
\node[pcm, above=6mm of saiTx] (pcmPb) {SAI\_Playback\\play 8ch};

% pipe zones
\node[pipe, fit=(b2a) (eqm) (b2b), label={[anchor=north west, font=\sffamily\bfseries\small, inner sep=2pt]north west:PIPE\_2}] {};
\node[pipe, fit=(pcmSrc) (b3a) (mbdrc) (b3b) (vol) (b3c) (drclim) (b3d), label={[anchor=north west, font=\sffamily\bfseries\small, inner sep=2pt]north west:PIPE\_3}] {};
\node[pipe, fit=(pcmTap), label={[anchor=north, font=\sffamily\bfseries\small, inner sep=1pt]north:PIPE\_5}] {};
\node[pipe, fit=(muxc) (b4) (saiTx), label={[anchor=north west, font=\sffamily\bfseries\small, inner sep=2pt]north west:PIPE\_4}] {};
\node[pipe, fit=(pcmPb), label={[anchor=south, font=\sffamily\bfseries\small, inner sep=1pt]south:PIPE\_6}] {};

% Arrows
\draw[arrmulti] (saiRx) -- (b1);
\draw[arrmulti] (b1) -- (pcmCap);

% SAI RX shared to PIPE_2
\draw[arrmulti] (saiRx.south) |- (b2a.west);

\draw[arr] (b2a) -- (eqm);
\draw[arr] (eqm) -- (b2b);
\draw[arrmulti] (b2b.east) -| node[above, font=\scriptsize, pos=0.1, align=center]{6ch monitor\\(slots 3-8)} ($(muxc.north west)+(0,2mm)$);

\draw[arr] (pcmSrc) -- (b3a);
\draw[arr] (b3a) -- (mbdrc);
\draw[arr] (mbdrc) -- (b3b);
\draw[arr] (b3b) -- (vol);
\draw[arr] (vol) -- (b3c);
\draw[arr] (b3c) -- (drclim);
\draw[arr] (drclim) -- (b3d);
\draw[arr] (b3d) -- node[above, font=\scriptsize]{2ch master} (muxc);

\draw[arr] (drclim.south) |- (pcmTap.west);

\draw[arr] (muxc) -- (b4);
\draw[arr] (b4) -- (saiTx);

\draw[arr] (pcmPb.south) -- node[right, font=\scriptsize]{override} (saiTx.north);

\end{tikzpicture}
\end{center}

\clearpage
\subsection{Détail de la chaîne master (PIPE\_3)}

\begin{center}
\begin{tikzpicture}[node distance=6mm]

\node[comp, fill=cdsp!30, minimum width=30mm] (mbdrc) {\textbf{multiband\_drc}\\3 bandes L-R linked\\bypass: \texttt{process\_enabled=0}\\$\Rightarrow$ \texttt{audio\_stream\_copy}};
\node[comp, fill=cdsp!30, minimum width=30mm, right=14mm of mbdrc] (vol) {\textbf{volume}\\DVC stereo\\bypass: \texttt{gain=0\,dB}\\$\Rightarrow$ mul par 1.0 fixed-pt};
\node[comp, fill=cdsp!30, minimum width=30mm, right=14mm of vol] (drc) {\textbf{drc}\\peak limiter \\bypass: \texttt{enabled=0}\\$\Rightarrow$ \texttt{drc\_default\_pass}};

\node[buf, left=6mm of mbdrc] (bi) {buf in\\2ch};
\node[buf, right=6mm of drc] (bo) {buf out\\2ch};

\draw[arr] (bi) -- (mbdrc);
\draw[arr] (mbdrc) -- node[above, font=\scriptsize]{2ch 32b} (vol);
\draw[arr] (vol) -- node[above, font=\scriptsize]{2ch 32b} (drc);
\draw[arr] (drc) -- (bo);

\node[below=1mm of mbdrc, font=\scriptsize\itshape, align=center]{crossover LR4 4 bandes\\+ DRC par bande (linked L-R)};
\node[below=1mm of vol, font=\scriptsize\itshape, align=center]{soft-ramp DVC,\\-$\infty$..+6 dB};
\node[below=1mm of drc, font=\scriptsize\itshape, align=center]{threshold +\\attack/release\\makeup gain};

\end{tikzpicture}
\end{center}

\textbf{Chaque composant a un flag de bypass bit-perfect} (prouvé par lecture du code):
\begin{itemize}[leftmargin=*, itemsep=2pt]
  \item \texttt{multiband\_drc\_default\_pass} dans \texttt{multiband\_drc\_generic.c:14} $=$ \texttt{audio\_stream\_copy(src,0,sink,0,ch*frames)}
  \item \texttt{drc\_default\_pass} dans \texttt{drc\_generic.c:473} $=$ idem
  \item \texttt{volume} à 0 dB (gain 0x40000000 en Q2.30) $=$ multiplication par 1.0, pas de loss
\end{itemize}

Le null-test T4 compare la sortie de \texttt{Master\_Tap} en mode bypass avec le signal envoyé à \texttt{Source\_Play}. Différence attendue : \textbf{0 dBFS (silence parfait)}. Si ce n'est pas le cas, il y a un bug dans le pipeline ou un composant a un side-effect non documenté.

\subsection{Détail du merge \texttt{mux} (PIPE\_4)}

Le composant \texttt{mux} (IPC3-compatible, voir \texttt{src/audio/mux/mux.c}) utilise une lookup table qui mappe des triplets \texttt{(stream\_id, in\_ch, out\_ch)} vers le sink. Config Phase 1a :

\begin{center}
\begin{tikzpicture}[node distance=3mm]

\node[buf] (m2ch) {2ch master};
\node[buf, below=10mm of m2ch] (m8ch) {6ch monitor\\(canaux 1-6 issus des mics 3-8 par ex.)};

\coordinate (midleft) at ($(m2ch)!0.5!(m8ch)$);
\node[comp, right=18mm of midleft, minimum width=22mm, minimum height=18mm, fill=cdsp!30, align=center] (mux) {\textbf{mux}\\lookup table};

\node[dai, right=14mm of mux, minimum height=22mm] (dai) {SAI TX 8ch\\\\slot 0: L master\\slot 1: R master\\slot 2: mon1\\slot 3: mon2\\\ldots\\slot 7: mon6};

\draw[arrmulti] (m2ch) -- (mux);
\draw[arrmulti] (m8ch) -- (mux);
\draw[arrmulti] (mux) -- (dai);

\node[right=14mm of dai, font=\sffamily\scriptsize, align=left]{
$\leftarrow$ TAC\#1 HP out L\\
$\leftarrow$ TAC\#1 HP out R\\
$\leftarrow$ TAC\#2 out 1\\
$\leftarrow$ TAC\#2 out 2\\
$\leftarrow$ TAC\#3 out 1\\
$\leftarrow$ TAC\#3 out 2\\
$\leftarrow$ TAC\#4 out 1\\
$\leftarrow$ TAC\#4 out 2};

\end{tikzpicture}
\end{center}

Lookup table (\texttt{sof\_mux\_config}) :
\begin{center}
\begin{tabular}{cccl}
\toprule
\texttt{stream\_id} & \texttt{in\_ch} & \texttt{out\_ch} & Commentaire \\
\midrule
0 & 0 & 0 & Master L $\rightarrow$ slot TDM 0 \\
0 & 1 & 1 & Master R $\rightarrow$ slot TDM 1 \\
1 & 0 & 2 & Monitor 1 $\rightarrow$ slot TDM 2 \\
1 & 1 & 3 & Monitor 2 $\rightarrow$ slot TDM 3 \\
1 & 2 & 4 & Monitor 3 $\rightarrow$ slot TDM 4 \\
1 & 3 & 5 & Monitor 4 $\rightarrow$ slot TDM 5 \\
1 & 4 & 6 & Monitor 5 $\rightarrow$ slot TDM 6 \\
1 & 5 & 7 & Monitor 6 $\rightarrow$ slot TDM 7 \\
\bottomrule
\end{tabular}
\end{center}

Stream 0 = sortie master chain (2ch). Stream 1 = 6 premiers canaux du monitor mic (on laisse tomber les 2 derniers mics sur ce merge, ou on garde les 8 via un mapping différent --- décision Phase 1a.3).

\clearpage
% =======================================================================
\section{Construction du fichier topology \texttt{sof-imx8mp-tac5212-masterlite.m4}}

\subsection{Principe}
La topology Phase 1a est décomposée en \textbf{4 pipelines} correspondant aux 4 PCMs. L'avantage : chaque pipeline est indépendant, on peut les activer/désactiver un par un pendant le développement.

\subsection{Mapping pipeline $\leftrightarrow$ PCM (V2)}

\begin{center}
\begin{tabularx}{\textwidth}{lllX}
\toprule
\textbf{Pipeline SOF} & \textbf{PCM exposé} & \textbf{PCM ID} & \textbf{Rôle} \\
\midrule
PIPE 1 + DAI\_ADD capture & \texttt{SAI\_Capture} & 0 & Dry mics 8ch (\texttt{PCM\_CAPTURE\_ADD}) \\
PIPE 3 (master chain in) & \texttt{Source\_Play} & 1 & Entrée stéréo master (\texttt{PCM\_PLAYBACK\_ADD}) \\
PIPE 5 (master chain out) & \texttt{Master\_Tap} & 2 & Sortie post-master (\texttt{PCM\_CAPTURE\_ADD}, co-scheduled avec PIPE 3) \\
PIPE 6 + DAI\_ADD playback & \texttt{SAI\_Playback} & 3 & Direct 8ch vers DAC (\texttt{PCM\_PLAYBACK\_ADD}) \\
\bottomrule
\end{tabularx}
\end{center}

\textit{\small Changement V2 : on ne peut PAS exposer SAI\_Capture et SAI\_Playback sous un même \texttt{PCM\_DUPLEX\_ADD} (un duplex = un seul nom pour les deux directions). V2 utilise 2 appels séparés, IDs 0 et 3.}

\subsection{Squelette m4 complet V2 (Phase 1a.1 MVP)}

Fichier : \texttt{sof/tools/topology/topology1/sof-imx8mp-tac5212-masterlite.m4}

Les changements V2 sont annotés avec \texttt{\# [V2]} dans le code.

\begin{lstlisting}[language={}, basicstyle=\ttfamily\scriptsize]
#
# Topology Phase 1a MVP V2 - 4 PCMs : SAI_Capture, SAI_Playback,
# Source_Play, Master_Tap + master chain stereo (mbdrc + volume + drc)
#
include(`utils.m4')
include(`dai.m4')
include(`pipeline.m4')
include(`sai.m4')
include(`pcm.m4')
include(`buffer.m4')
include(`common/tlv.m4')
include(`sof/tokens.m4')
include(`platform/imx/imx8.m4')

# ---------------------------------------------------------------------
# PIPE 1 : Dry mics capture (8ch)
# ---------------------------------------------------------------------
PIPELINE_PCM_ADD(sof/pipe-passthrough-capture.m4,
    1, 0, 8, s32le,
    2000, 0, 0,
    48000, 48000, 48000)

# ---------------------------------------------------------------------
# PIPE 6 : Direct playback (8ch) - PCM ID 3 pour eviter collision
# ---------------------------------------------------------------------
PIPELINE_PCM_ADD(sof/pipe-passthrough-playback.m4,
    6, 3, 8, s32le,
    2000, 0, 0,
    48000, 48000, 48000)

# ---------------------------------------------------------------------
# PIPE 3 : Master chain input side (Source_Play PCM 1)
# pipeline CUSTOM local : HOST -> mbdrc -> volume -> drc -> buf
# ---------------------------------------------------------------------
PIPELINE_PCM_ADD(sof-imx8mp-tac5212-pipe-masterchain.m4,
    3, 1, 2, s32le,
    2000, 0, 0,
    48000, 48000, 48000)

# ---------------------------------------------------------------------
# PIPE 5 : Master chain output side (Master_Tap PCM 2)
# [V2] co-scheduled avec PIPE 3 via PIPELINE_SCHED_COMP_3
#      -> partage le scheduler, pas de race entre timers independants
# ---------------------------------------------------------------------
PIPELINE_PCM_ADD(sof-imx8mp-tac5212-pipe-mastertap.m4,
    5, 2, 2, s32le,
    2000, 0, 0,
    48000, 48000, 48000,
    SCHEDULE_TIME_DOMAIN_TIMER,
    PIPELINE_SCHED_COMP_3)

# ---------------------------------------------------------------------
# DAI_ADD : SAI7 RX + TX (identique baseline drc2ms qui fonctionne)
# ---------------------------------------------------------------------
DAI_ADD(sof/pipe-dai-capture.m4,
    1, SAI, 7, tac5212-hifi,
    PIPELINE_SINK_1, 2, s32le,
    2000, 0, 0, SCHEDULE_TIME_DOMAIN_DMA)

DAI_ADD(sof/pipe-dai-playback.m4,
    6, SAI, 7, tac5212-hifi,
    PIPELINE_SOURCE_6, 2, s32le,
    2000, 0, 0, SCHEDULE_TIME_DOMAIN_DMA)

# ---------------------------------------------------------------------
# PCM exports V2 : 4 noms distincts, IDs 0/1/2/3
# [V2] remplace PCM_DUPLEX_ADD par 2 ajouts separes
# ---------------------------------------------------------------------
PCM_CAPTURE_ADD (SAI_Capture,  0, PIPELINE_PCM_1)
PCM_PLAYBACK_ADD(Source_Play,  1, PIPELINE_PCM_3)
PCM_CAPTURE_ADD (Master_Tap,   2, PIPELINE_PCM_5)
PCM_PLAYBACK_ADD(SAI_Playback, 3, PIPELINE_PCM_6)

# ---------------------------------------------------------------------
# Graph : connecter le buffer sink de PIPE 3 a la source de PIPE 5
# ---------------------------------------------------------------------
SectionGraph."master-chain-to-tap" {
    index "0"
    lines [
        `dapm(PIPELINE_SINK_5, PIPELINE_SOURCE_3)'
    ]
}

# ---------------------------------------------------------------------
# DAI config SAI7 TDM 8 slots x 32 bits @ 48kHz (identique drc2ms.m4)
# ---------------------------------------------------------------------
DAI_CONFIG(SAI, 7, 0, tac5212-hifi,
    SAI_CONFIG(DSP_A, SAI_CLOCK(mclk, 12288000, codec_mclk_in),
        SAI_CLOCK(bclk, 12288000, codec_consumer),
        SAI_CLOCK(fsync, 48000, codec_consumer),
        SAI_TDM(8, 32, 255, 255),
        SAI_CONFIG_DATA(SAI, 7, 0)))
\end{lstlisting}

\subsection{Pipelines m4 custom à créer (Phase 1a.1)}

Deux pipelines custom n'existent pas en upstream SOF et doivent être écrits dans \texttt{meta-local/recipes-kernel/linux/files/} (OU dans \texttt{sof/tools/topology/topology1/sof/} si acceptable). Règle : on ne touche PAS les upstream \texttt{sof/pipe-*.m4}.

\paragraph{\texttt{sof-imx8mp-tac5212-pipe-masterchain.m4} V2}
Pipeline d'entrée de la chaîne master. Valeurs concrètes : 2 canaux, s32le, 96 frames/period (2 ms @ 48 kHz).
\begin{lstlisting}[language={}, basicstyle=\ttfamily\scriptsize]
# Master chain input V2: host PCM -> mbdrc -> volume -> drc -> buf
# PIPELINE_CHANNELS = 2, PIPELINE_FORMAT = s32le, SCHEDULE_PERIOD = 2000us
include(`utils.m4')
include(`buffer.m4')
include(`pcm.m4')
include(`pipeline.m4')
include(`bytecontrol.m4')
include(`multiband_drc.m4')
include(`volume.m4')
include(`drc.m4')

# Host PCM playback endpoint (2ch, 3 periodes de sink)
W_PCM_PLAYBACK(PCM_ID, Source Play, 2, 0, SCHEDULE_CORE)

# --- multiband_drc ---
define(MBDRC_priv, concat(`mbdrc_bytes_', PIPELINE_ID))
define(MBDRC_CTRL, concat(`mbdrc_control_', PIPELINE_ID))
include(`mbdrc_coef_default.m4')
C_CONTROLBYTES(MBDRC_CTRL, PIPELINE_ID,
    CONTROLBYTES_OPS(bytes, 258 binds the control to bytes get/put, 258, 258),
    CONTROLBYTES_EXTOPS(258 binds the control to bytes get/put, 258, 258),
    , , ,
    CONTROLBYTES_MAX(, 1024),
    ,
    MBDRC_priv)
W_MULTIBAND_DRC(0, PIPELINE_FORMAT, 2, 2, SCHEDULE_CORE,
    LIST(`   ', "MBDRC_CTRL"))

# --- volume (DVC master) ---
# [V2] ramp NONE pour bit-perfect bypass
W_VOLUME(0, PIPELINE_FORMAT, 2, 2, SCHEDULE_CORE, 0, SOF_VOLUME_RAMP_NONE)

# --- drc (final limiter) ---
define(DRC_priv, concat(`drc_bytes_', PIPELINE_ID))
define(DRC_CTRL, concat(`drc_control_', PIPELINE_ID))
include(`drc_coef_limiter_default.m4')
C_CONTROLBYTES(DRC_CTRL, PIPELINE_ID,
    CONTROLBYTES_OPS(bytes, 258 binds the control to bytes get/put, 258, 258),
    CONTROLBYTES_EXTOPS(258 binds the control to bytes get/put, 258, 258),
    , , ,
    CONTROLBYTES_MAX(, 1024),
    ,
    DRC_priv)
W_DRC(0, PIPELINE_FORMAT, 2, 2, SCHEDULE_CORE,
    LIST(`   ', "DRC_CTRL"))

# --- Buffers concretes ---
# Un buffer par lien : host | mbdrc-out | vol-out | drc-out(shared)
W_BUFFER(0, COMP_BUFFER_SIZE(3,
    COMP_SAMPLE_SIZE(PIPELINE_FORMAT), PIPELINE_CHANNELS,
    COMP_PERIOD_FRAMES(PCM_MAX_RATE, SCHEDULE_PERIOD)),
    PLATFORM_HOST_MEM_CAP)
W_BUFFER(1, COMP_BUFFER_SIZE(2,
    COMP_SAMPLE_SIZE(PIPELINE_FORMAT), PIPELINE_CHANNELS,
    COMP_PERIOD_FRAMES(PCM_MAX_RATE, SCHEDULE_PERIOD)),
    PLATFORM_COMP_MEM_CAP)
W_BUFFER(2, COMP_BUFFER_SIZE(2,
    COMP_SAMPLE_SIZE(PIPELINE_FORMAT), PIPELINE_CHANNELS,
    COMP_PERIOD_FRAMES(PCM_MAX_RATE, SCHEDULE_PERIOD)),
    PLATFORM_COMP_MEM_CAP)
W_BUFFER(3, COMP_BUFFER_SIZE(2,
    COMP_SAMPLE_SIZE(PIPELINE_FORMAT), PIPELINE_CHANNELS,
    COMP_PERIOD_FRAMES(PCM_MAX_RATE, SCHEDULE_PERIOD)),
    PLATFORM_COMP_MEM_CAP)

# Graph : host -> B0 -> mbdrc -> B1 -> volume -> B2 -> drc -> B3
P_GRAPH(pipe-master-chain, PIPELINE_ID,
    LIST(`  ',
        `dapm(N_BUFFER(0), N_PCMP(PCM_ID))',
        `dapm(N_MULTIBAND_DRC(0), N_BUFFER(0))',
        `dapm(N_BUFFER(1), N_MULTIBAND_DRC(0))',
        `dapm(N_VOLUME(0), N_BUFFER(1))',
        `dapm(N_BUFFER(2), N_VOLUME(0))',
        `dapm(N_DRC(0), N_BUFFER(2))',
        `dapm(N_BUFFER(3), N_DRC(0))'))

# Export buffer 3 comme SOURCE de la pipeline pour le graph externe
indir(`define', concat(`PIPELINE_SOURCE_', PIPELINE_ID), N_BUFFER(3))
indir(`define', concat(`PIPELINE_PCM_', PIPELINE_ID), Source Play PCM_ID)

# PCM capabilities : 2ch s32le @ 48k, 192..16384 bytes period, 65536 buf
PCM_CAPABILITIES(Source Play PCM_ID, CAPABILITY_FORMAT_NAME(PIPELINE_FORMAT),
    PCM_MIN_RATE, PCM_MAX_RATE, 2, 2, 2, 16, 192, 16384, 65536, 65536)

undefine(`MBDRC_CTRL')
undefine(`MBDRC_priv')
undefine(`DRC_CTRL')
undefine(`DRC_priv')
\end{lstlisting}

\paragraph{\texttt{sof-imx8mp-tac5212-pipe-mastertap.m4} V2}
Pipeline de sortie côté Master\_Tap, simple capture depuis un buffer externe (partagé avec pipe-masterchain buf\#3 via DAPM).
\begin{lstlisting}[language={}, basicstyle=\ttfamily\scriptsize]
# Master tap V2: buf (shared) -> host PCM Master_Tap
# Co-scheduled avec PIPE 3 via PIPELINE_SCHED_COMP_3 passe dans
# PIPELINE_PCM_ADD (cf. sof-imx8mp-tac5212-masterlite.m4)
include(`utils.m4')
include(`buffer.m4')
include(`pcm.m4')
include(`pipeline.m4')

# Host PCM capture endpoint (2ch, 0 sink periodes, 2 source periodes)
W_PCM_CAPTURE(PCM_ID, Master Tap, 0, 2, SCHEDULE_CORE)

# Buffer d'entree (sera lie au N_BUFFER(3) de masterchain via DAPM externe)
W_BUFFER(0, COMP_BUFFER_SIZE(2,
    COMP_SAMPLE_SIZE(PIPELINE_FORMAT), PIPELINE_CHANNELS,
    COMP_PERIOD_FRAMES(PCM_MAX_RATE, SCHEDULE_PERIOD)),
    PLATFORM_HOST_MEM_CAP)

# Graph : buf -> host
P_GRAPH(pipe-master-tap, PIPELINE_ID,
    LIST(`   ',
        `dapm(N_PCMC(PCM_ID), N_BUFFER(0))'))

# Export buffer 0 comme SINK de la pipeline pour le graph externe
indir(`define', concat(`PIPELINE_SINK_', PIPELINE_ID), N_BUFFER(0))
indir(`define', concat(`PIPELINE_PCM_', PIPELINE_ID), Master Tap PCM_ID)

PCM_CAPABILITIES(Master Tap PCM_ID, CAPABILITY_FORMAT_NAME(PIPELINE_FORMAT),
    PCM_MIN_RATE, PCM_MAX_RATE, 2, 2, 2, 16, 192, 16384, 65536, 65536)
\end{lstlisting}

\subsection{Évolutions 1a.2 et 1a.3}

\paragraph{1a.2 --- Ajout PIPE 2 (monitor direct mic $\rightarrow$ HP)}
\begin{itemize}[leftmargin=*, itemsep=1pt]
  \item Nouveau pipeline \texttt{sof-imx8mp-tac5212-pipe-monitor.m4} : SAI RX $\rightarrow$ eq\_iir$\times$8 $\rightarrow$ buf partagé avec PIPE 4
  \item Attach au scheduling de PIPE 1 via \texttt{PIPELINE\_SCHED\_COMP\_1} (pattern kwd)
  \item Ajouter dans le graph : \texttt{dapm(PIPELINE\_SINK\_2, PIPELINE\_SOURCE\_1)} pour partager SAI RX
\end{itemize}

\paragraph{1a.3 --- Ajout PIPE 4 (mux merge master + monitor)}
\begin{itemize}[leftmargin=*, itemsep=1pt]
  \item Nouveau pipeline \texttt{sof-imx8mp-tac5212-pipe-muxout.m4} : 2ch master + 6ch monitor $\rightarrow$ \texttt{mux} $\rightarrow$ SAI TX
  \item Remplace ou coexiste avec PIPE 6 (SAI\_Playback direct)
  \item Lookup table \texttt{mux} configurée via bytes control (cf. section 3.6 de ce document)
\end{itemize}

\clearpage
% =======================================================================
\section{Budget latence}

\subsection{Chemins audio et latences attendues}

\begin{center}
\begin{tabular}{lrl}
\toprule
\textbf{Chemin} & \textbf{Latence} & \textbf{Contrainte} \\
\midrule
Mic $\rightarrow$ HP direct (PIPE\_2 $\rightarrow$ mux $\rightarrow$ SAI TX) & \textbf{4{,}0~ms} & $\leq$ 4~ms \ding{51} \\
Mic $\rightarrow$ \texttt{SAI\_Capture} host (PIPE\_1) & 4{,}0~ms & N/A \\
\texttt{Source\_Play} $\rightarrow$ HP (PIPE\_3 $\rightarrow$ PIPE\_4) & 8{,}0~ms & libre (playback mastering) \\
\texttt{Source\_Play} $\rightarrow$ \texttt{Master\_Tap} (PIPE\_3 $\rightarrow$ PIPE\_5) & 8{,}0~ms & tolérable pour NPU (update 10 Hz) \\
Mic $\rightarrow$ \texttt{Master\_Tap} (si matrice Phase 1b) & 6--8~ms & tolérable \\
\bottomrule
\end{tabular}
\end{center}

\subsection{Décomposition mic\,$\rightarrow$\,HP direct (objectif 4 ms)}

\begin{center}
\begin{tikzpicture}[node distance=2mm]
\node[draw, fill=cdsp!15, minimum width=20mm, minimum height=8mm, font=\sffamily\footnotesize, align=center] (s1) {TAC5212 ADC};
\node[draw, fill=chw!15, right=1mm of s1, minimum width=18mm, minimum height=8mm, font=\sffamily\footnotesize, align=center] (s2) {DMA SAI RX};
\node[draw, fill=cdsp!15, right=1mm of s2, minimum width=28mm, minimum height=8mm, font=\sffamily\footnotesize, align=center] (s3) {Pipeline 2~ms\\(eq\_iir + mux)};
\node[draw, fill=chw!15, right=1mm of s3, minimum width=18mm, minimum height=8mm, font=\sffamily\footnotesize, align=center] (s4) {DMA SAI TX};
\node[draw, fill=cdsp!15, right=1mm of s4, minimum width=20mm, minimum height=8mm, font=\sffamily\footnotesize, align=center] (s5) {TAC5212 DAC};

\node[below=1mm of s1, font=\sffamily\scriptsize] {$<$100~\textmu s};
\node[below=1mm of s2, font=\sffamily\scriptsize] {1 burst 64 smp};
\node[below=1mm of s3, font=\sffamily\scriptsize] {1 cycle $=$ 2~ms};
\node[below=1mm of s4, font=\sffamily\scriptsize] {1 burst 64 smp};
\node[below=1mm of s5, font=\sffamily\scriptsize] {$<$100~\textmu s};

\node[below=8mm of s3, font=\sffamily] {\textbf{Total $\approx$ 4~ms} (2 cycles DSP period de 2~ms)};
\end{tikzpicture}
\end{center}

\textbf{Marge de manœuvre} : aucune. Si on ajoute un composant dans PIPE\_2 (ex. compresseur per-ch), on dépasse 4~ms. Le fine-EQ par canal (3 biquads IIR) est acceptable car négligeable en temps de calcul ($\ll$ 100~\textmu s par canal).

Si T8 échoue (mic$\rightarrow$HP $>$ 4~ms mesuré), mitigation :
\begin{itemize}[leftmargin=*, itemsep=1pt]
  \item passer la DSP period à 1~ms $\rightarrow$ 2~ms total, mais charge CPU doublée
  \item retirer eq\_iir de PIPE\_2 et le faire uniquement dans PIPE\_1 (capture) $\rightarrow$ perte d'EQ sur monitoring
  \item laisser TAC5212 faire toute l'EQ fine (3 biquads par canal en hard) et supprimer eq\_iir du DSP
\end{itemize}

\clearpage
% =======================================================================
\section{Tests T1--T9 --- commandes complètes}

Scripts et binaires stockés dans \texttt{/root/tests/} sur la board (convention projet). WAVs de référence générés avec \texttt{sox} sur la board.

\subsection{T0 --- Pré-validation (V2)}
\textbf{But} : vérifier que les prérequis board + outils sont en place avant T1.

\begin{lstlisting}[language=bash]
# 1. Nom exact de la carte ALSA (peut differer de softac5212tdm)
cat /proc/asound/cards

# 2. Enumere les PCMs exposes par la topology
arecord -l 2>&1 | grep -A1 tac5212
aplay   -l 2>&1 | grep -A1 tac5212
# Doivent apparaitre: SAI_Capture (cap), Source_Play (play),
#                      Master_Tap (cap), SAI_Playback (play)

# 3. Controls ALSA attendus
amixer -c <CARD> controls | grep -iE 'master|bypass|volume|mbdrc|limiter'

# 4. Outils de test
which sox arecord aplay amixer sof-ctl sof-logger python3
test -f /lib/firmware/imx/sof/sof-imx8m.ldc || \
    echo "WARN: sof-imx8m.ldc manquant, T9 ne pourra pas lire les traces"
\end{lstlisting}

\textbf{Pass si} : 4 PCMs trouvés, tous les outils disponibles, \texttt{.ldc} présent.
\textbf{Si nom carte différent}, remplacer \texttt{softac5212tdm} par le nom réel dans la suite des tests, ou exporter \texttt{CARD=<nom>} et utiliser \texttt{plughw:\$CARD,X} partout.

\subsection{T1 --- Capture dry 8ch}
\textbf{But} : vérifier que \texttt{SAI\_Capture} expose bien les 8 canaux mics bruts, sans glitch.
\begin{lstlisting}[language=bash]
# Génère un signal sur IN1 (ch1) via phone/PC branché
arecord -D plughw:softac5212tdm,SAI_Capture -c 8 -r 48000 -f S32_LE \
        -d 5 /root/tests/out/dry8.wav
sox /root/tests/out/dry8.wav -n stat   # RMS / peak par canal
\end{lstlisting}
\textbf{Pass si} : ch1 a du signal ($>$ -40~dBFS), ch2--8 au noise floor ($<$ -100~dBFS), 0 xrun dans \texttt{dmesg}.

\subsection{T2 --- Playback master}
\textbf{But} : jouer un tone stéréo à travers la chaîne master.
\begin{lstlisting}[language=bash]
sox -n -r 48000 -c 2 -b 32 /root/tests/in/tone440.wav synth 3 sine 440 \
    vol -10dB
aplay -D plughw:softac5212tdm,Source_Play /root/tests/in/tone440.wav
\end{lstlisting}
\textbf{Pass si} : tone audible sur TAC\#1 HP OUT, pas de clic, pas d'xrun.

\subsection{T3 --- Master tap $=$ master chain output}
\textbf{But} : vérifier que \texttt{Master\_Tap} capte bien le signal processé.
\begin{lstlisting}[language=bash]
arecord -D plughw:softac5212tdm,Master_Tap -c 2 -r 48000 -f S32_LE \
        -d 5 /root/tests/out/tap.wav &
sleep 1
aplay -D plughw:softac5212tdm,Source_Play /root/tests/in/tone440.wav
wait
sox /root/tests/out/tap.wav -n stat
\end{lstlisting}
\textbf{Pass si} : tap.wav contient le tone 440 Hz à niveau cohérent avec la chaîne (vol 0 dB + mbdrc \& drc configurés).

\subsection{T4 --- Null test bypass (V2, aligné cross-correlation)}
\textbf{But} : prouver que la chaîne master en bypass est bit-perfect.

\textbf{[V2]} Le pipeline master a une latence fixe ($\approx$ 6-8 ms en bypass). Un \texttt{sox -m} direct entre \texttt{tone440.wav} source et la capture échoue car les signaux ne sont pas alignés temporellement --- on mesure le décalage, pas la bit-perfection. V2 utilise un script Python avec cross-correlation pour aligner au sample près.

\begin{lstlisting}[language=bash]
# 1. Activer bypass + attendre ramp volume (100ms)
amixer -c softac5212tdm cset name='Master MBDRC Bypass' on
amixer -c softac5212tdm cset name='Master Volume' 0dB
amixer -c softac5212tdm cset name='Master Limiter Bypass' on
sleep 0.2     # [V2] laisser converger soft-ramp volume

# 2. Stimulus : impulsion courte + silence (aide alignement)
sox -n -r 48000 -c 2 -b 32 /root/tests/in/stim.wav \
    synth 0.001 square 1000 vol -6dB pad 0.1 3

# 3. Capture parallele du Master_Tap
arecord -D plughw:softac5212tdm,Master_Tap -c 2 -r 48000 -f S32_LE \
        -d 4 /root/tests/out/bypass.wav &
PID=$!
sleep 0.2
aplay -D plughw:softac5212tdm,Source_Play /root/tests/in/stim.wav
wait $PID

# 4. Null-test aligne
python3 /root/tests/null_test.py \
        /root/tests/in/stim.wav \
        /root/tests/out/bypass.wav \
        --threshold -140
\end{lstlisting}

\textbf{Script Python} \texttt{/root/tests/null\_test.py} :
\begin{lstlisting}[language={}, basicstyle=\ttfamily\scriptsize]
#!/usr/bin/env python3
"""Null test aligne par cross-correlation. Exit 0 si RMS < threshold."""
import sys, argparse, numpy as np, soundfile as sf
from scipy.signal import correlate

ap = argparse.ArgumentParser()
ap.add_argument('source'); ap.add_argument('capture')
ap.add_argument('--threshold', type=float, default=-140.0)
ap.add_argument('--channel', type=int, default=0)
args = ap.parse_args()

src, sr = sf.read(args.source, dtype='int32')
cap, sr = sf.read(args.capture, dtype='int32')
if src.ndim == 2: src = src[:, args.channel]
if cap.ndim == 2: cap = cap[:, args.channel]

# Cross-correlation pour trouver le delay capture vs source
# (on prend un segment court pour eviter ram explosion)
N = min(48000, len(src), len(cap))
xc = correlate(cap[:N].astype(np.float64), src[:N].astype(np.float64), mode='full')
lag = np.argmax(np.abs(xc)) - (N - 1)
print(f"Detected lag: {lag} samples ({lag/48:.2f} ms)")

# Align et tronque a la plus petite taille commune
if lag > 0:
    cap_a = cap[lag:]
    src_a = src[:len(cap_a)]
else:
    src_a = src[-lag:]
    cap_a = cap[:len(src_a)]
L = min(len(src_a), len(cap_a))
diff = src_a[:L].astype(np.float64) - cap_a[:L].astype(np.float64)

rms = np.sqrt(np.mean(diff**2))
peak = np.max(np.abs(diff))
db_rms = 20 * np.log10(rms / 2**31 + 1e-30)
db_peak = 20 * np.log10(peak / 2**31 + 1e-30)
print(f"Null RMS : {db_rms:.1f} dBFS | Peak : {db_peak:.1f} dBFS")

sys.exit(0 if db_rms < args.threshold else 1)
\end{lstlisting}

\textbf{Pass si} : \texttt{null\_test.py} exit 0 (RMS $<$ -140 dBFS). \textbf{Tolerance dégradée} : si la SIMD HiFi4 introduit du rounding LSB, accepter $<$ -120 dBFS (1 LSB sur 32 bits $\approx$ -186 dBFS). Si Fail $>$ -80 dBFS : vrai bug à investiguer.

\subsection{T5 --- Duplex stress 30 min}
\textbf{But} : stabilité sur long cours.
\begin{lstlisting}[language=bash]
# Fichier de 30 min stéréo
sox -n -r 48000 -c 2 -b 32 /root/tests/in/stress.wav synth 1800 pinknoise vol -20dB
arecord -D plughw:softac5212tdm,SAI_Capture -c 8 -r 48000 -f S32_LE \
        -d 1800 /root/tests/out/stress_cap.wav &
arecord -D plughw:softac5212tdm,Master_Tap -c 2 -r 48000 -f S32_LE \
        -d 1800 /root/tests/out/stress_tap.wav &
aplay -D plughw:softac5212tdm,Source_Play /root/tests/in/stress.wav &
wait
dmesg | grep -i xrun | wc -l   # doit retourner 0
\end{lstlisting}
\textbf{Pass si} : 0 xrun sur 30 min, \texttt{stat} des 3 fichiers cohérent.

\subsection{T6 --- ALSA controls live}
\textbf{But} : modifier volume pendant playback sans glitch.
\begin{lstlisting}[language=bash]
aplay -D plughw:softac5212tdm,Source_Play /root/tests/in/tone440.wav &
sleep 1
amixer -c softac5212tdm cset name='Master Volume' -20dB
sleep 1
amixer -c softac5212tdm cset name='Master Volume' 0dB
wait
\end{lstlisting}
\textbf{Pass si} : volume chute de 20 dB à l'oreille, remonte, pas de clic.

\subsection{T7 --- Bytes control mbdrc}
\textbf{But} : pousser une config \texttt{multiband\_drc} différente.
\begin{lstlisting}[language=bash]
# Preset "heavy compression" préparé
sof-ctl -D softac5212tdm -n 'Master MBDRC Config' -b -s /root/tests/presets/mbdrc_heavy.bin

arecord -D plughw:softac5212tdm,Master_Tap -c 2 -r 48000 -f S32_LE \
        -d 5 /root/tests/out/tap_heavy.wav &
sleep 1
aplay -D plughw:softac5212tdm,Source_Play /root/tests/in/tone440.wav
wait

# Comparer dynamique
sox /root/tests/out/tap_heavy.wav -n stat | grep 'Maximum'
\end{lstlisting}
\textbf{Pass si} : peak du WAV "heavy" $<$ peak "default" (compression plus forte).

\subsection{T8 --- Latence mic$\rightarrow$HP $\leq$ 4 ms}
\textbf{But} : mesurer physiquement la latence.
\begin{lstlisting}[language=bash]
# Cable loopback : OUT1 (TAC#1 HP) --> IN2 (TAC#1 mic 2)
# Envoie un click (impulsion) via Source_Play, capture via SAI_Capture ch2

sox -n -r 48000 -c 2 -b 32 /root/tests/in/click.wav synth 0.001 square 1000 \
    vol -6dB pad 1 0

arecord -D plughw:softac5212tdm,SAI_Capture -c 8 -r 48000 -f S32_LE \
        -d 3 /root/tests/out/loop.wav &
sleep 0.1
aplay -D plughw:softac5212tdm,Source_Play /root/tests/in/click.wav
wait

# Calcule le delay par cross-correlation
python3 /root/tests/measure_latency.py /root/tests/in/click.wav \
        /root/tests/out/loop.wav --channel 1
\end{lstlisting}
\textbf{Pass si} : script Python reporte delay $\leq$ 6~ms \textbf{[V2]} (cible théorique 4 ms + jitter timer SOF $\pm$500 µs + group delay SigmaDelta TAC5212 $\approx$ 1.5 ms).

\textbf{[V2 notes]} Le 4 ms strict n'est atteignable que si le TAC5212 est en mode ultra-low-latency filter (group delay $\approx$ 2.8 samples). En mode linear-phase par défaut (16 samples group delay), on vise 5-6 ms.

\subsection{T9 --- Charge CPU DSP}
\textbf{But} : mesurer l'occupation CPU HiFi4 en fonctionnement nominal.
\begin{lstlisting}[language=bash]
# Avec SOF-logger, lire les traces contenant CPC (cycles per chunk)
sof-logger -f /lib/firmware/imx/sof/sof-imx8m.ldc \
           -o /tmp/sof.log -r 5 &
LOGPID=$!

# Charger le DSP : aplay + arecord + arecord master
aplay -D plughw:softac5212tdm,Source_Play /root/tests/in/stress.wav &
arecord -D plughw:softac5212tdm,SAI_Capture -c 8 -r 48000 -f S32_LE \
        -d 10 /dev/null &
arecord -D plughw:softac5212tdm,Master_Tap -c 2 -r 48000 -f S32_LE \
        -d 10 /dev/null &
wait

kill $LOGPID
grep CPC /tmp/sof.log | awk '{print $NF}' | sort -u | tail
\end{lstlisting}
\textbf{Pass si} : utilisation rapportée $<$ \textbf{35\%} de la capacité HiFi4 \textbf{[V2]} (relevé de 30\% pour compter l'overhead des memcpy host bridges i.MX8MP, même avec AP2AP actif).

\textbf{Prérequis} : \texttt{sof-tools} doit être dans \texttt{IMAGE\_INSTALL} (voir §6 livrables) et le fichier \texttt{.ldc} doit être déployé à côté du \texttt{.ri} dans \texttt{/lib/firmware/imx/sof/}.

\subsection{Script unique \texttt{test-phase1a.sh}}
Un script orchestre T1--T9 séquentiellement, produit un rapport Markdown. Exit 0 si tous passent.

\clearpage
% =======================================================================
\section{Livrables --- fichiers et leur localisation}

\subsection{Arborescence des fichiers Phase 1a}

\begin{lstlisting}[language=bash]
/home/michael/yocto-nxp-debix/
+-- sof/                                    # sous-module SOF, branche feature
|   +-- app/boards/imx8mp_evk_mimx8ml8_adsp.conf   # override Kconfig si besoin
|   +-- tools/topology/topology1/
|       +-- sof-imx8mp-tac5212-masterlite.m4       # LA topology Phase 1a
|
+-- meta-local/
|   +-- recipes-kernel/linux/files/
|   |   +-- tac5212.c                       # driver codec existant
|   |   +-- tac5212-register-map.h          # map des registres (Phase 1a)
|   +-- recipes-support/
|   |   +-- tac5212-init/
|   |   |   +-- tac5212-init_1.0.bb
|   |   |   +-- files/
|   |   |       +-- tac5212-init.c          # service boot I2C
|   |   |       +-- tac5212-init.service    # systemd unit
|   |   |       +-- default.conf            # preset biquad+AGC+DVC
|   |   +-- sof-presets/
|   |   |   +-- sof-presets_1.0.bb
|   |   |   +-- files/
|   |   |       +-- mbdrc_default.bin       # preset mbdrc par defaut
|   |   |       +-- mbdrc_heavy.bin         # preset mbdrc compression forte
|   |   |       +-- drc_limiter_default.bin # preset final limiter
|   |   +-- phase1a-tests/
|   |       +-- phase1a-tests_1.0.bb
|   |       +-- files/
|   |           +-- test-phase1a.sh         # orchestrateur T1-T9
|   |           +-- measure_latency.py      # calcul delay par xcorr
|   |           +-- t1_capture_dry.sh       # T1 standalone
|   |           +-- t2_playback_master.sh
|   |           +-- ...
|   +-- recipes-core/images/
|       +-- imx-image-full.bbappend         # IMAGE_INSTALL += new packages
|
+-- PHASE_1A_SPEC.pdf                       # CE DOCUMENT
+-- PHASE_1A_SPEC.tex
+-- MASTERING_PLATFORM_PLAN.md              # plan vivant
+-- DSP_ARCHITECTURES.pdf                   # comparatif A/B/C
\end{lstlisting}

\subsection{Fichiers SOF à créer / modifier}

\begin{center}
\begin{tabularx}{\textwidth}{lX}
\toprule
\textbf{Fichier} & \textbf{Action} \\
\midrule
\texttt{sof/tools/topology/topology1/sof-imx8mp-tac5212-masterlite.m4} & CRÉER (nouvelle topology Phase 1a) \\
\texttt{sof/app/boards/imx8mp\_evk\_mimx8ml8\_adsp.conf} & MODIFIER (ajouter \texttt{CONFIG\_COMP\_MULTIBAND\_DRC=y} et \texttt{CONFIG\_COMP\_MUX=y} si pas déjà présent ; laisser IPC3 par défaut) \\
\bottomrule
\end{tabularx}
\end{center}

\textbf{Rappel règle projet} : ne JAMAIS modifier les fichiers upstream \texttt{sof/tools/topology/topology1/sof/*.m4}. Si on a besoin d'un pipeline custom inconnu en amont, il vit dans \texttt{meta-local/recipes-kernel/linux/files/} sous un nom unique.

\subsection{Fichiers Yocto à créer}

\begin{center}
\begin{tabularx}{\textwidth}{lX}
\toprule
\textbf{Fichier} & \textbf{Rôle} \\
\midrule
\texttt{meta-local/recipes-support/tac5212-init/tac5212-init\_1.0.bb} & Recette paquet service I\textsuperscript{2}C \\
\texttt{meta-local/recipes-support/tac5212-init/files/tac5212-init.c} & Service C --- écrit les biquads/AGC/DVC au boot \\
\texttt{meta-local/recipes-support/tac5212-init/files/tac5212-init.service} & systemd unit (Type=oneshot, avant sof-firmware) \\
\texttt{meta-local/recipes-support/tac5212-init/files/default.conf} & Preset biquad/AGC/DVC au format texte \\
\texttt{meta-local/recipes-support/sof-presets/files/mbdrc\_default.bin} & Bytes blob mbdrc (format SOF IPC3) \\
\texttt{meta-local/recipes-support/sof-presets/files/mbdrc\_heavy.bin} & Preset compression forte (pour T7) \\
\texttt{meta-local/recipes-support/sof-presets/files/drc\_limiter\_default.bin} & Preset limiter -0{,}1 dBFS ceiling \\
\texttt{meta-local/recipes-support/phase1a-tests/files/test-phase1a.sh} & Orchestrateur de tests \\
\texttt{meta-local/recipes-support/phase1a-tests/files/measure\_latency.py} & Script Python pour T8 (cross-correlation) \\
\texttt{meta-local/recipes-core/images/imx-image-full.bbappend} & \textbf{[V2]} \texttt{IMAGE\_INSTALL += tac5212-init sof-presets phase1a-tests \textbf{sof-tools python3-numpy python3-scipy python3-soundfile}} \\
\bottomrule
\end{tabularx}
\end{center}

\subsection{Déploiement du \texttt{.ldc} (V2)}
Le fichier \texttt{zephyr.ri.ldc} (log dictionary) produit par le build SOF doit être déployé sur la board pour permettre à \texttt{sof-logger} de lire les traces (requis par T9). Commande à ajouter après chaque rebuild firmware :
\begin{lstlisting}[language=bash]
scp /tmp/build-imx8m/zephyr/zephyr.ri.ldc \
    root@192.168.0.9:/lib/firmware/imx/sof/sof-imx8m.ldc
\end{lstlisting}
Sans \texttt{.ldc}, \texttt{sof-logger} produit un stream binaire illisible.

\clearpage
% =======================================================================
\section{Commandes de build et déploiement}

\subsection{Build topology uniquement (cas courant Phase 1a)}

Quand on modifie juste le \texttt{.m4} de la topology, on ne reconstruit pas le firmware SOF complet --- seulement le \texttt{.tplg}.

\begin{lstlisting}[language=bash]
cd /home/michael/yocto-nxp-debix/sof/tools/topology/topology1

# 1. Preprocess m4
m4 -I . -I m4 -I common -I platform/common -I sof \
    sof-imx8mp-tac5212-masterlite.m4 \
    > /tmp/sof-imx8mp-tac5212-masterlite.conf

# 2. Compile conf -> tplg
alsatplg -c /tmp/sof-imx8mp-tac5212-masterlite.conf \
         -o /tmp/sof-imx8mp-tac5212-masterlite.tplg

# 3. Deploy (le kernel charge /lib/firmware/imx/sof-tplg/sof-imx8mp-tac5212.tplg)
scp /tmp/sof-imx8mp-tac5212-masterlite.tplg \
    root@192.168.0.9:/lib/firmware/imx/sof-tplg/sof-imx8mp-tac5212.tplg

# 4. Reboot (chargement tplg se fait au probe driver SOF)
ssh root@192.168.0.9 systemctl reboot

# 5. Verifier
ssh root@192.168.0.9 "dmesg | grep -i 'sof\|tplg' | tail -20"
\end{lstlisting}

\textbf{Piège sed} : si on crée une variante par sed \texttt{'s/8000/4000/g'}, attention, le pattern matche aussi \texttt{48000} et \texttt{12288000}. Toujours utiliser \texttt{Edit} avec un pattern précis.

\subsection{Build firmware SOF complet (cas où on modifie Kconfig ou sources)}

Si on ajoute \texttt{CONFIG\_COMP\_MULTIBAND\_DRC} dans le board config, il faut rebuilder le firmware \texttt{.ri} :

\begin{lstlisting}[language=bash]
cd /home/michael/yocto-nxp-debix

# 1. Build firmware (Zephyr west)
west build -d /tmp/build-imx8m -b imx8mp_evk/mimx8ml8/adsp sof/app

# 2. Sign avec rimage
/tmp/build-rimage/rimage -o /tmp/build-imx8m/zephyr/zephyr.ri -e \
    -k sof/keys/otc_private_key.pem \
    -c sof/tools/rimage/config/imx8m.toml \
    /tmp/build-imx8m/zephyr/zephyr.elf

# 3. CRITICAL: concat xman + ri (sinon IPC 108/20 error)
cat /tmp/build-imx8m/zephyr/zephyr.ri.xman \
    /tmp/build-imx8m/zephyr/zephyr.ri \
    > /tmp/build-imx8m/zephyr/zephyr.ri.full

# 4. Deploy (fichier cible: sof-imx8m.ri, SANS 'p')
scp /tmp/build-imx8m/zephyr/zephyr.ri.full \
    root@192.168.0.9:/lib/firmware/imx/sof/sof-imx8m.ri

# 5. Reboot
ssh root@192.168.0.9 systemctl reboot

# 6. Verifier hash change
ssh root@192.168.0.9 "dmesg | grep 'Firmware info'"
\end{lstlisting}

\subsection{Build complet Yocto (intégration dans l'image)}

Pour les paquets Yocto (service \texttt{tac5212-init}, presets, tests) :

\begin{lstlisting}[language=bash]
cd /home/michael/yocto-nxp-debix

# Initialiser environnement (une fois par session)
EULA=1 DISTRO=fsl-imx-xwayland MACHINE=imx8mpevk \
    source imx-setup-release.sh -b Model_AB_Infinity

# Rebuild d'un paquet uniquement
bitbake tac5212-init
bitbake phase1a-tests
bitbake sof-presets

# Full image (pour deploy complet SD)
bitbake imx-image-full

# Rebuild complet si recette cassée
bitbake -c cleansstate <recipe>   # JAMAIS cleanall (cf. memoire projet)
bitbake <recipe>
\end{lstlisting}

\subsection{Deploy paquet seul sur la board}

\begin{lstlisting}[language=bash]
# Trouver le .deb produit
find Model_AB_Infinity/tmp/deploy/deb/imx8mpevk -name 'tac5212-init_*.deb'

# scp + install sur la board
scp Model_AB_Infinity/tmp/deploy/deb/imx8mpevk/tac5212-init_*.deb \
    root@192.168.0.9:/tmp/
ssh root@192.168.0.9 "dpkg -i /tmp/tac5212-init_*.deb && \
                      systemctl daemon-reload && \
                      systemctl restart tac5212-init"
\end{lstlisting}

\subsection{Deploy full image (rare, SD card)}

\begin{lstlisting}[language=bash]
# Flash SD
dd if=Model_AB_Infinity/tmp/deploy/images/imx8mpevk/imx-image-full-imx8mpevk.wic \
   of=/dev/sdX bs=4M status=progress conv=fsync
sync
\end{lstlisting}

\clearpage
% =======================================================================
\section{Protocole d'exécution Phase 1a}

\subsection{Règles durables}
\begin{enumerate}[leftmargin=*, itemsep=2pt]
  \item \textbf{Aucune modification de source SOF sans \texttt{critic\_analyze} préalable.}
  \item \textbf{Un changement = un commit = une validation utilisateur.}
  \item \textbf{Backup systématique} avant toute édition d'un fichier existant (\texttt{cp fichier.m4 /tmp/fichier.m4.bak}).
  \item \textbf{Jamais toucher \texttt{sof/tools/topology/topology1/sof/*.m4}} (fichiers upstream SOF). Nos pipelines custom vont dans \texttt{meta-local/} ou dans un fichier topology unique \texttt{sof-imx8mp-tac5212-masterlite.m4}.
  \item \textbf{Tests stockés dans \texttt{/root/tests/}} sur la board, \textbf{jamais dans \texttt{/tmp/}} (mémoire projet).
\end{enumerate}

\subsection{Séquence 1a.1 --- MVP (V2)}
\begin{enumerate}[leftmargin=*, itemsep=2pt]
  \item Design \texttt{sof-imx8mp-tac5212-masterlite.m4} minimal (PIPE\_1 + PIPE\_3 + PIPE\_5 + PIPE\_6).
  \item \texttt{critic\_analyze} sur le diff topology.
  \item \textbf{[V2]} \textbf{Gate compile-test} sur machine dev :
    \begin{lstlisting}[language=bash]
cd sof/tools/topology/topology1
# copier les pipelines custom (masterchain + mastertap) dans topology1/
m4 -I . -I m4 -I common -I platform/common -I sof \
    sof-imx8mp-tac5212-masterlite.m4 > /tmp/test.conf \
  && alsatplg -c /tmp/test.conf -o /tmp/test.tplg \
  && echo "COMPILE OK, $(stat -c%s /tmp/test.tplg) bytes"
\end{lstlisting}
    \textbf{STOP} si erreur m4 ou alsatplg. Corriger le squelette AVANT deploy.
  \item \textbf{[V2]} T0 sur la board : vérifier nom carte, PCMs, outils.
  \item Deploy tplg, reboot.
  \item Exécuter T1 (capture dry).
  \item Exécuter T2 (playback master avec preset par défaut).
  \item Exécuter T3 (master tap reçoit bien le signal).
  \item Exécuter T4 (null test bypass aligné).
  \item Exécuter T6 (alsactrl live).
  \item Exécuter T7 (bytes control mbdrc).
  \item Validation utilisateur $\rightarrow$ commit MVP.
\end{enumerate}

\subsection{Séquence 1a.2 --- Monitor direct}
\begin{enumerate}[leftmargin=*, itemsep=2pt]
  \item Ajouter PIPE\_2 (SAI RX shared $\rightarrow$ eq\_iir$\times$8 $\rightarrow$ buf partagé avec PIPE\_4).
  \item Pattern DAPM cross-pipeline (\texttt{PIPELINE\_SCHED\_COMP\_1}) comme dans kwd.
  \item \texttt{critic\_analyze} sur le diff.
  \item Build + deploy + reboot.
  \item Exécuter T8 (latence mic$\rightarrow$HP).
  \item Validation utilisateur.
\end{enumerate}

\subsection{Séquence 1a.3 --- Merge mux}
\begin{enumerate}[leftmargin=*, itemsep=2pt]
  \item Remplacer PIPE\_6 direct-playback par PIPE\_4 avec \texttt{mux} (2ch master + 6ch monitor).
  \item Préparer la lookup table \texttt{mux}.
  \item \texttt{critic\_analyze}.
  \item Build + deploy + reboot.
  \item Exécuter T5 (stress 30 min) + T9 (CPU).
  \item Validation utilisateur $\rightarrow$ Phase 1a OK, passage à Phase 1b.
\end{enumerate}

\subsection{Critères de sortie Phase 1a}
\begin{itemize}[leftmargin=*, itemsep=1pt]
  \item T1--T9 tous passent ;
  \item \texttt{test-phase1a.sh} exit 0 ;
  \item 0 xrun sur un run de 30 min ;
  \item Latence mic$\rightarrow$HP $\leq$ 4 ms mesurée (T8) ;
  \item CPU DSP $<$ 30\% (T9) ;
  \item Documentation à jour : \texttt{MASTERING\_PLATFORM\_PLAN.md} marque Phase 1a comme \emph{done}.
\end{itemize}

\clearpage
% =======================================================================
\section{Vue courte des phases suivantes}

\subsection{Phase 1b --- Matrice complète}
Remplacer le routing fixe Phase 1a par une vraie matrice N\,$\times$\,M via composants \texttt{mixer} parallèles (un par bus de sortie) $+$ gains configurables via controls ALSA.

Pré-requis : valider Phase 1a stable + Phase 2 USB opérationnelle (car la matrice a plus de sens quand il y a plus d'inputs).

\subsection{Phase 2 --- USB gadgets}
2 gadgets UAC2 via \texttt{configfs} sur les deux contrôleurs USB du i.MX8MP :
\begin{itemize}[leftmargin=*, itemsep=1pt]
  \item USB1 : 8in/8out pour DAW ASIO
  \item USB2 : 2in/2out pour iOS
\end{itemize}
Bridges ALSA entre les gadgets et le DSP via \texttt{alsaloop} (ou bridge C custom si latence serrée).

\subsection{Phase 3 --- Effets send (A53 JACK)}
Serveur JACK2 avec 4 plugins LV2 (reverb, chorus, flanger, delay). Bridges ALSA $\leftrightarrow$ DSP via 4 PCMs capture + 4 PCMs playback supplémentaires.

Sélection des plugins à faire parmi :
\begin{itemize}[leftmargin=*, itemsep=1pt]
  \item Reverb : \texttt{dragonfly-reverb}, \texttt{calf-reverb}
  \item Chorus : \texttt{calf-multichorus}
  \item Flanger : \texttt{calf-flanger}
  \item Delay : \texttt{x42-delay}, \texttt{calf-delay}
\end{itemize}

\subsection{Phase 4 --- NPU mastering}
Training d'un modèle TFLite (from scratch) pour le mastering assisté :
\begin{itemize}[leftmargin=*, itemsep=1pt]
  \item Input : FFT du master (2048 pts, 20 ms)
  \item Output : gains EQ 8-10 bandes + seuils mbdrc par bande + amount exciter
  \item Update rate : 10 Hz
\end{itemize}

\subsection{Phase 5 --- Harmonizer + polish}
\begin{itemize}[leftmargin=*, itemsep=1pt]
  \item Harmonizer pitch-shift : plugin LV2 \texttt{rubberband} chaîné en A53
  \item Exciter : plugin LV2 ou module SOF custom (à décider sur CPU budget)
  \item Flutter UI avancée : presets, scenes, métering en temps réel
\end{itemize}

\clearpage
% =======================================================================
\section{Annexe --- composants SOF mobilisés}

\subsection{Liste et rôles}
\begin{center}
\begin{tabularx}{\textwidth}{llX}
\toprule
\textbf{Composant} & \textbf{Kconfig} & \textbf{Rôle Phase 1a} \\
\midrule
\texttt{dai-copier} & natif & Interface SAI7 RX/TX $\leftrightarrow$ buffers DSP \\
\texttt{host-copier} & natif & Interface HOST PCM $\leftrightarrow$ buffers DSP \\
\texttt{eq\_iir} & \texttt{COMP\_IIR} & 3 biquads IIR/ch (fine EQ, flat par défaut) \\
\texttt{mux} & \texttt{COMP\_MUX} & Merge 2ch master + 6ch monitor $\rightarrow$ 8ch SAI TX \\
\texttt{multiband\_drc} & \texttt{COMP\_MULTIBAND\_DRC} & Compresseur multibande stéréo (bypass bit-perfect dispo) \\
\texttt{volume} & natif & Digital Volume Control master (0 dB bit-perfect) \\
\texttt{drc} & \texttt{COMP\_DRC} & Final limiter (bypass bit-perfect dispo) \\
\bottomrule
\end{tabularx}
\end{center}

\subsection{Bypass bit-perfect --- preuves code}
\begin{center}
\begin{tabularx}{\textwidth}{lll}
\toprule
\textbf{Composant} & \textbf{Fichier} & \textbf{Fonction bypass} \\
\midrule
\texttt{multiband\_drc} & \texttt{multiband\_drc\_generic.c:14} & \texttt{multiband\_drc\_default\_pass} ($\equiv$ \texttt{audio\_stream\_copy}) \\
\texttt{drc} & \texttt{drc\_generic.c:473} & \texttt{drc\_default\_pass} ($\equiv$ \texttt{audio\_stream\_copy}) \\
\texttt{volume} & \texttt{volume.c} & gain Q2.30 $=$ \texttt{0x40000000} $=$ 1{,}0 $\Rightarrow$ multiplication par 1 \\
\bottomrule
\end{tabularx}
\end{center}

\subsection{ALSA controls exposés Phase 1a}

\begin{center}
\begin{tabularx}{\textwidth}{llX}
\toprule
\textbf{Nom} & \textbf{Type} & \textbf{Plage / effet} \\
\midrule
\texttt{SAI EQ Ch<N> Coeffs} & bytes blob & 3 biquads IIR (coefficients normalisés) \\
\texttt{Master Volume} & integer TLV dB & -$\infty$..+6 dB, 0.5 dB step \\
\texttt{Master MBDRC Config} & bytes blob & Structure \texttt{sof\_multiband\_drc\_config} \\
\texttt{Master MBDRC Bypass} & switch 0/1 & Active \texttt{multiband\_drc\_default\_pass} \\
\texttt{Master Limiter Config} & bytes blob & Structure \texttt{sof\_drc\_config} \\
\texttt{Master Limiter Bypass} & switch 0/1 & Active \texttt{drc\_default\_pass} \\
\bottomrule
\end{tabularx}
\end{center}

\subsection{Références datasheet TAC5212}
Voir \texttt{tac5212.pdf} dans la racine projet. Sections clés :
\begin{itemize}[leftmargin=*, itemsep=1pt]
  \item §7.3.9.1 ADC signal chain (biquads, AGC, HPF, gain cal)
  \item §7.3.9.2 DAC signal chain (biquads, DRC DAC, volume, interpolation)
  \item §8.2 Programmable coefficient registers
\end{itemize}

\vfill
\begin{center}
\textit{Document de référence Phase 1a. Tout écart à cette spec doit être tracé dans \texttt{MASTERING\_PLATFORM\_PLAN.md} et re-soumis à \texttt{critic\_analyze}.}
\end{center}

\end{document}
